% -----------------------------------------------------------
\section{Urim and Thummim}
% -----------------------------------------------------------
	\label{URIMTHUMMIM}
	
% % ----------------------
% \subsection{See also}
% % ----------------------

% % ----------------------
% \subsection{1828 American Dictionary of the English Language} % *1828
% % ----------------------
	% \label{URIMTHUMMIM-1828}

	% \begin{compactitem}
		% \item
	% \end{compactitem}

% % ----------------------
% \subsection{Oxford English Dictionary} % *OED
% % ----------------------
	% \label{URIMTHUMMIM-OED}

	% \begin{compactitem}
		% \item[]
	% \end{compactitem}
	
% % ----------------------
% \subsection{Restoration Lexicon} % *RL *RESTORATION LEXICON
% % ----------------------
	% \label{URIMTHUMMIM-OED}

	% \begin{compactitem}
		% \item[]
	% \end{compactitem}

% ----------------------
\subsection{Scriptural Usage} % *SCRIPTURES
% ----------------------
	\label{URIMTHUMMIM-SCRIPTURES}

	% -----
	\subsubsection{Old Testament} % *OT
	% -----
		% \label{URIMTHUMMIM-OT}
	
		% % --
		% \paragraph{KJV}
		% % --
			% \label{URIMTHUMMIM-OT-KJV}
	
			% \begin{tabular}{ll}
				% Total occurrences in the OT & 
			% \end{tabular}
			
			% \begin{compactdesc}
				% \item[]
			% \end{compactdesc}
			
		% --
		\paragraph{Hebrew Bible}
		% --
			\label{URIMTHUMMIM-HEBREW}
		
			%
			\subparagraph{\texorpdfstring{\he{'UriyM}}{}} % אוּרִים
			%
				\label{URIM} % whatever is in step bible without periods
			
				\begin{compactdesc}
					\item[HALOT 25a, PDF 82] \textbf{} % Use the 1994 PDF
						\begin{compactitem}
							\item \textbf{Urim}
								\begin{compactitem}
									\item The entry mentions, apparently as possibilities, ``instrument for casting lots...Arabic, blunt arrows...various coloured stones.''
								\end{compactitem}
						\end{compactitem}
					\item[BDB 22a, PDF 41] \textbf{}
						\begin{compactitem}
							\item \textbf{Urim} (plural intensifier, mostly with article \textit{the Urim}, and mostly joined with \he{t*um*iyM}...\he{hA'UriyM w:hat*um*iyM} were put into the \he{.ho+sEN} was a little bag of pouch worn on breast of high priest, to hold the \textit{Urim \& Thummim}; the name \he{.ho+sEN ham*i+s:p*A.t} was given because of decisions made by that which was within it; thus, Eleazar was to inquire of \he{y} for Joshua \he{b*:mi+s:p*a.t hA'UriyM} (\hyperref[NUMBERS-27-21]{Numbers 27:21})
								\begin{compactitem}
									\item The entry also notes the LXX of \hyperref[DEUTERONOMY-33-8]{Deuteronomy 33:8}, \gr{καὶ τῷ Λευι εἶπεν Δότε λευι δήλους αὐτοῦ Καὶ ἀλήθειαν αὐτοῦ τῷ ἀνδρὶ τῷ ὁσίῳ}, and the Vulgate, \textit{Levi quoque ait perfectio tua et doctrina tua viro sancto tuo quem probasti in Temptatione}, but in the HALOT entry it suggests these sources are just guessing at the meaning.
									\item The \he{.ho+sEN} is the ``breastpiece'': ``The breast piece o fAaron's priestly wardrobe is described in great detail in the tabernacle texts of Exodus (see esp. \hyperref[EXODUS-28-5]{28:5--30} and \hyperref[EXODUS-39-8]{39:8--21}) and is also mentioned in one passage in Leviticus (\hyperref[LEVITICUS-8-8]{8:8}) dealing with the tabernacle. The breastpiece was made of the same kind of material as the ephod: gold, blue, purple, and scarlet woolen threads interwoven with fine linen. The workmanship was of a certain skillful type (\textit{hoseb}) also used for the ephod. The woven fabric formed a double piece of material, a span 9 or 10 inches square when folded over. The reason for the doubled fabric was that the breastpiece, in addition to its symbolic ritual value, also served as a container for the Urim and Thummim'' (ABD 1182).
								\end{compactitem}
						\end{compactitem}
					\item[DCH PDF 165] \textbf{} % don't include the actual page number, they repeat from volume to volume
						\begin{compactitem}
							\item \textbf{Urim}...sacred oracle in high priest's breastplate, usually + \he{t*um*iyM} \textit{Thummim}
						\end{compactitem}
				\end{compactdesc}
				
				% \begin{compactdesc}
					% \item[Some OT uses] \textbf{}
						% \begin{compactdesc}
							% \item[]
						% \end{compactdesc}
				% \end{compactdesc}
				
				% %
				% \subparagraph{TDNT: \texorpdfstring{\gr{sample word}}{}  (3:536--556)}
				% %
					% \label{URIM-ABD}
					
			%
			\subparagraph{\texorpdfstring{\he{t*um*iyM}}{}} % תֻּמִּים
			%
				\label{TUMMIM} % whatever is in step bible without periods
				
				The best I can gather from these entries is that \he{t*um*iyM} comes from \he{tmM}, ``to be complete'' (Qal), ``to complete, perfect, accomplish, carry out'' (Hiphil).
			
				\begin{compactdesc}
					\item[HALOT 1750b, PDF 1806] \textbf{} % Use the 1994 PDF
						\begin{compactitem}
							\item plural from \he{t*oM} (KBL) or an independent tantum plurale...epithet of \he{'UriyM}...(\hyperref[NUMBERS-27-21]{Numbers 27:21} and \hyperref[1SAMUEL-28-6]{1 Samuel 28:6} \he{'UriyM} stands figuratively as a part for the whole for \he{'UriyM w:tum*iyM}..the etymological sense of \he{t*um*iyM} is uncertain; the proposal that \he{'UriyM} belongs to the root \he{'rr} and correspondingly that \he{t*u} to the root \he{tmM} leaves open the possibility of considering the word-pair cursed + innocent...for another view see e.g. Eissfeldt...who takes the sense of Urim and Thummim (without explanation) as a clarification of guilt and confirmation of innocence. In practice \he{t*um*iyM} functioned together with \he{'UriyM} to serve as a means of ascertaining through an oracle by lot the divine will using arrows, or rather differently marked stones
						\end{compactitem}
					\item[BDB 1070b, PDF 1089] \textbf{}
						\begin{compactitem}
							\item \he{t*um*iyM} in \he{'UriyM w:tum*iyM} (meaning dubious, see \he{'wriyM})
								\begin{compactitem}
									\item BDB lists this word as \#4 in \he{t*oM}, ``\textbf{completeness, integrity}.'' 
								\end{compactitem}
						\end{compactitem}
					\item[DCH PDF 5484a] \textbf{} % don't include the actual page number, they repeat from volume to volume
						\begin{compactitem}
							\item \textit{Thummim}
								\begin{compactitem}
									\item DCH lists this word under \he{tmM}, ``\textbf{to be complete},'' but the word only appears at the very end of the entry in a list of derivatives from \he{tmM} that includes \he{t*oM} and other words.
								\end{compactitem}
						\end{compactitem}
				\end{compactdesc}
			
		% % --
		% \paragraph{Septuagint}
		% % --
			% \label{URIMTHUMMIM-LXX}
		
			% %
			% \subparagraph{\texorpdfstring{\gr{sample word}}{}}
			% %
		
	% % -----
	% \subsubsection{New Testament} % *NT
	% % -----
		% \label{URIMTHUMMIM-NT}
	
		% % --
		% \paragraph{KJV}
		% % --
			% \label{URIMTHUMMIM-NT-KJV}		
	
			% \begin{tabular}{ll}
				% Total occurrences in the NT & 
			% \end{tabular}
			
			% \begin{compactdesc}
				% \item[]
			% \end{compactdesc}
			
		% % --
		% \paragraph{Greek New Testament} % *GREEK
		% % --
			% \label{URIMTHUMMIM-GREEK}
		
			% %
			% \subparagraph{\texorpdfstring{\gr{sample word}}{}}
			% %
				% \label{KATALLAGE}
			
				% \begin{compactdesc}
					% \item[BDAG ] \textbf{}
						% \begin{compactitem}
							% \item 
						% \end{compactitem}
					% \item[CGL , PDF ] \textbf{}
						% \begin{compactitem}
							% \item 
						% \end{compactitem}
					% \item[LSJ , PDF ] \textbf{}
						% \begin{compactitem}
							% \item 
						% \end{compactitem}
				% \end{compactdesc}
				
				% \begin{tabular}{ll}
					% Total occurrences in the NT & 
				% \end{tabular}
				
				% \begin{compactdesc}
					% \item[Some NT uses] \textbf{}
						% \begin{compactdesc}
							% \item[]
						% \end{compactdesc}
				% \end{compactdesc}
				
			% %
			% \subparagraph{TDNT: \texorpdfstring{\gr{sample word}}{}  (3:536--556)}
			% %
				% \label{KATALLAGE-TDNT}
		
	% % -----
	% \subsubsection{Book of Mormon} % *BOM
	% % -----
		% \label{URIMTHUMMIM-BOM}
	
		% \begin{tabular}{ll}
			% Total occurrences in the BOM & 
		% \end{tabular}
		
		% \begin{compactdesc}
			% \item[]
		% \end{compactdesc}
		
	% % -----
	% \subsubsection{Doctrine and Covenants} % *DC
	% % -----
		% \label{URIMTHUMMIM-DC}
	
		% \begin{tabular}{ll}
			% Total occurrences in the DC & 
		% \end{tabular}
		
		% \begin{compactdesc}
			% \item[]
		% \end{compactdesc}
		
	% % -----
	% \subsubsection{Pearl of Great Price} % *PGP
	% % -----
		% \label{URIMTHUMMIM-PGP}
	
		% \begin{tabular}{ll}
			% Total occurrences in the PGP & 
		% \end{tabular}
		
		% \begin{compactdesc}
			% \item[]
		% \end{compactdesc}
		
% % ----------------------
% \subsection{Key Phrases} % *KEYPHRASES *PHRASES
% % ----------------------
	% \label{URIMTHUMMIM-PHRASES}

	% \begin{compactdesc}
		% \item[] \textbf{\#times} | list
			% % \begin{compactdesc}
				% % \item[Bible anchor verses] \phantom{.}
					% % \begin{compactdesc}
						% % \item[Romans 5:5] \gr{}
					% % \end{compactdesc}
				% % \item[Commentary] ---
			% % \end{compactdesc}
	% \end{compactdesc}
	
% % ----------------------
% \subsection{Bible Dictionaries} % *DICTIONARIES
% % ----------------------
	% \label{URIMTHUMMIM-BD}

% % ----------------------
% \subsection{Restoration Usage} % *RESTORATION
% % ----------------------
	% \label{URIMTHUMMIM-RESTORATION}

	% % -----
	% \subsubsection{Joseph Smith} % *JS
	% % -----
		% \label{URIMTHUMMIM-JS}
	
		% \begin{compactdesc}
			% \item[{\hyperref[KEY]{Date/Source}}]
		% \end{compactdesc}
		
	% % -----
	% \subsubsection{Temple Ordinances}
	% % -----
		% \label{URIMTHUMMIM-TEMPLE}
	
		% \begin{compactdesc}
			% \item[{\hyperref[KEY]{Date/Source}}]
		% \end{compactdesc}
	
	% % -----
	% \subsubsection{Brigham Young} % *BY
	% % -----
		% \label{URIMTHUMMIM-BY}
	
		% \begin{compactdesc}
			% \item[{\hyperref[KEY]{Date/Source}}]
		% \end{compactdesc}
	
% ----------------------
\subsection{Commentary and Studies} % *COMMENTARY *STUDIES
% ----------------------
	\label{URIMTHUMMIM-COMMENTARY}

	% % -----
	% \subsubsection{Key Points}
	% % -----
		% \label{URIMTHUMMIM-KEYS}
	
		% \begin{compactitem}
			% \item 
		% \end{compactitem}
		
	% -----
	\subsubsection{What was the Urim and Thummim?}
	% -----
		\label{URIMTHUMMIM-what}
		
		Ann Taves, \textit{Revelatory Events} (2016, Princeton UP), 37--38:
		
		\begin{quote}
			Instead of viewing spectacles, directors, interpreters, and Urim and Thummim as designating \textit{different} objects, we can view them as pro gressively elaborated interpretations of the stone(s) that Smith used in the translation process.6 Each of the terms contextualized the stone(s) in progressively larger narratives that simultaneously undercut competing treasure-seeking narratives and gave the stone(s) new meaning and significance within a developing Mormon narrative. We can place these terms and their associated stories in a plausible developmental sequence. The first story was most likely Joseph's story of recovering something he could use to translate the plates, whether ``spectacles'' or ``a glass he obtained with the plates.'' The story of an object recovered along with or after the plates was most likely followed by the Book of Mormon stories; first the stories in Mosiah about the ``interpreters''---``two stones which was fastened into the two rims of a bow'' that could be used to translate ancient records, and then the stories in Alma about the ``directors''---``a stone'' that ``Gazelem'' would use to bring forth the content of the Jaredites' twenty- four plates. In 1833 W. W. Phelps connected the objects described in the Book of Mormon with the biblical ``Urim and Thummim'' (e.g., 1 Sam. 14:41--42). Smith subsequently embraced this idea in his canonized history, as did the later tradition, such that ``Urim and Thummim'' were adopted as generic terms designating objects that were used to discern the mind or will of the deity (Van Wagoner and Walker, 49--50, 53). Later revelations further integrated these stories by interpreting Gazelem as a reference to Joseph Smith (D\&C 78:9; 82:11; 104:26, 43--46) and elaborating on the cosmic significance of the Urim and Thummim (D\&C 130).
		\end{quote}
		
	% -----
	\subsubsection{Restoration usage of the term ``Urim and Thummim''}
	% -----
		\label{URIMTHUMMIM-restoration}
		
		\begin{compactitem}
			\item A footnote from JSP D1:84: ``In this version of the revelation, the use of `Urim and Thummim' (rather than the Book of Mormon term `interpreters' or the term `spectacles,' which JS used in 1829 and 1832) is probably a later redaction since `Urim and Thummim' does not appear in JS's writings before 1833. The revisions in this section may in part be correcting errors made while copying from a source text that had itself been revised.''
			\item As far as I know---and this comes from text-searching the full transcript---the term ``Urim and Thummim'' does not appear in the 1833 BC.
		\end{compactitem}